% BAN TIENG VIET cua 07_system_model.tex (2026-10-08). So lieu giu nguyen.
\section{Mô hình hệ thống}
\label{sec:model}

\subsection{Cân bằng nước của ruộng}

Lượng nước trong tầng rễ $w_t$ [mm] biến đổi theo giờ dưới mô hình thùng chứa
FAO-56~\cite{fao56}:
\begin{equation}
w_{t+1}=\mathrm{clip}_{\le FC}\!\left(w_t+P_t-RO_t-ET_{c,t}+u_t\right),
\label{eq:bucket}
\end{equation}
với lượng bốc thoát hơi nước của cây trồng và hệ số ứng suất nước
\begin{equation}
ET_{c,t}=K_{s,t}\,K_c\,ET_{0,t},
\qquad
K_{s,t}=\frac{FC-w_t}{(1-p)\,\mathrm{TAW}},
\label{eq:etc}
\end{equation}
trong đó $ET_{0,t}$ là bốc thoát hơi nước tham chiếu theo giờ
(Penman-Monteith), $P_t$ là lượng mưa, $RO_t$ là dòng chảy mặt (phương pháp
đường cong số SCS, tính theo ngày rồi chia đều cho các giờ của một trận mưa),
và $u_t$ là độ sâu tưới áp dụng trong giờ $t$. Lượng nước vượt quá sức chứa ẩm
đồng ruộng $FC$ bị toán tử cắt bỏ loại đi, nên bảo toàn khối lượng đúng tuyệt
đối và được kiểm tra như một bất biến. Độ cạn kiệt của tầng rễ là
\begin{equation}
d_t=\frac{FC-w_t}{FC-WP},\qquad \mathrm{span}=FC-WP,
\label{eq:depl}
\end{equation}
với $WP$ là điểm héo; đối với loại đất thịt pha sét limon tham chiếu của đồng
bằng, $\mathrm{TAW}=FC-WP=70$~mm, $FC=160$~mm, $WP=90$~mm.

\subsection{Phát lại một nhật ký khai báo}
\label{sec:replay}

Giả sử một nhật ký khai báo cung cấp số hạng tưới của \eqref{eq:bucket}.
\emph{Toán tử phát lại} $\mathcal{R}$ ánh xạ một nhật ký sang quỹ đạo mà nó
hàm ý, bằng cách lặp \eqref{eq:bucket} dưới dữ liệu khí tượng quan trắc
$(P_t,ET_{0,t})$:
\begin{equation}
\mathcal{R}(L)=\{(\hat w_t,\hat d_t)\}_{t=0}^{T},
\label{eq:replay}
\end{equation}
với $\hat w_0=w_{\mathrm{init}}$ và
\begin{equation}
u_t=\sum_{(t_i,u_i)\in L}u_i\,\mathbb{1}[t_i=t],
\label{eq:uterm}
\end{equation}
sao cho $\hat w_t$ là lượng nước tầng rễ mà nhật ký \emph{khai} rằng ruộng đã
có, và $\hat d_t$ là độ cạn kiệt của nó theo \eqref{eq:depl}; dấu mũ đánh dấu
các đại lượng được tái dựng từ nhật ký chứ không đo được. Việc phát lại là
tất định và chỉ dùng dữ liệu tái phân tích công khai, nên hai kiểm toán viên
phát lại cùng một nhật ký sẽ thu được cùng một quỹ đạo và phán quyết có thể
tái lập. Dấu mũ trừ chỉ trạng thái ngay \emph{trước} khi khai báo của giờ đó
được áp dụng, như trong $\hat w_{t}^{-}$: lượng nước sẵn có để tiếp nhận một
lần tưới, và đó là thứ bộ lọc phải xem xét.

Toán tử này cũng là thứ làm cho phương án trở nên rẻ: không cảm biến nào bị
truy vấn, và đầu vào duy nhất là nhật ký, chuỗi thời tiết công khai cùng một
hồ sơ đất. Điểm yếu của nó, được trình bày ở Mục~\ref{sec:timeshift}, là nó
chỉ kiểm tra được một quỹ đạo có \emph{khả dĩ} hay không, chứ không bao giờ
biết được nó có \emph{đã xảy ra} hay không.
